\documentclass[11pt]{article}
\usepackage{ideastyle}

\begin{document}

\ideatitle{GuideVoice}{A conversational walking guide for blind and low-vision users.}

\section{The Idea}
GPS apps tell you \emph{``turn left in 200 feet''} but not that there's a
construction barrier, which door is the entrance, or where the crosswalk button
is. GuideVoice combines turn-by-turn directions with the phone camera: point it
ahead and ask \emph{``What's in front of me?''} or \emph{``Where's the door?''}, and
Grok Voice answers in a calm, conversational way. This directly answers an opening
ceremony prompt: \emph{``What types of audio or haptic navigation would blind or
low-vision users benefit from?''}

\section{Track Fit}
\begin{tabular}{@{}P{0.28\textwidth}P{0.66\textwidth}@{}}
\toprule
\req{Built with Cursor}{Whole app built in Cursor.}
\req{Grok Voice / Imagine}{\textbf{Voice} is the entire interface.}
\req{Space data}{Not used. Allowed by the rules, but the track is pitched around space, so judges may weigh this.}
\req{Navigation theme}{Direct fit, and an accessibility prompt from the slides.}
\req{Also eligible}{Big Red, Software, Design, People's Choice.}
\bottomrule
\end{tabular}

\section{Features}
\subsection{MVP}
\begin{itemize}
  \item Voice-only start: \emph{``Take me to Klarman Hall.''}
  \item Walking directions spoken at each turn.
  \item ``What's ahead?'' button: snap a photo, get a short spoken description focused on obstacles, doors, and signs.
  \item Large-text, high-contrast, screen-reader-friendly UI.
\end{itemize}
\subsection{Stretch}
\begin{itemize}
  \item Haptic patterns for left/right/arrived (phone vibration).
  \item Continuous mode that describes changes every few seconds.
  \item Read signs and room numbers aloud.
\end{itemize}

\section{Tech Stack and Data}
\begin{itemize}
  \item Frontend: mobile web app (PWA) with camera and geolocation.
  \item Routing: Google Maps Directions API (workshop Saturday 1:00 PM) or OpenStreetMap routing.
  \item Vision: an image-understanding model (Grok's vision-capable chat model, to verify).
  \item Voice: Grok Voice API.
\end{itemize}

\section{Limitations and Risks}
\begin{itemize}
\limitation{Safety}{Wrong or late descriptions could lead a blind user into danger. It must never replace a cane, guide dog, or orientation training.}{Position it as an assistant, include a clear disclaimer, and never say ``it's safe to cross.''}
\limitation{Latency}{Photo upload, AI description, and speech take seconds; a walking person moves several meters meanwhile.}{Keep descriptions short; use the on-demand button instead of continuous mode for the MVP.}
\limitation{GPS accuracy}{Phone GPS is off by 5--20 meters near buildings, which matters a lot to someone who cannot see.}{Use the camera to confirm landmarks (``the entrance is on your right'').}
\limitation{No real users to test with}{Without blind or low-vision testers, design choices are guesses.}{Follow established accessibility guidelines; the Saturday 3:00 PM Assistive Tech workshop may help.}
\limitation{Existing products}{Be My AI, Seeing AI, and Google Lookout already describe scenes.}{Differentiate by combining scene description \emph{with} route guidance in one conversation.}
\limitation{Grok API unknowns}{Image-input support and Voice API details are not verified.}{Confirm at the Friday Grok/Cursor workshop.}
\end{itemize}

\section{Scorecard}
\begin{tabular}{@{}P{0.25\textwidth}P{0.08\textwidth}P{0.6\textwidth}@{}}
\toprule
\rating{Demo-able}{5}{Works live, indoors, at the judging table.}
\rating{Buildable in 24h}{4}{Mostly API integration.}
\rating{Navigation theme}{5}{Direct accessibility prompt.}
\rating{Wow factor}{4}{Judges can close their eyes and try it.}
\rating{Technical risk (5 = riskiest)}{2}{Low to moderate.}
\bottomrule
\end{tabular}

\section{Verdict}
\textbf{Top non-space pick.} Meaningful, demos live indoors, and matches a prompt
from the opening slides. Be careful with safety claims.

\end{document}
